\documentclass[10pt,a4paper]{amsart}
\usepackage{amsmath,amssymb,amsthm,graphicx}
\usepackage{verbatim}
\newenvironment{reference}{\comment}{\endcomment}
%\newenvironment{reference}{}{}
\newenvironment{slogan}{\comment}{\endcomment}
\newenvironment{history}{\comment}{\endcomment}

% For commutative diagrams we use Xy-pic
\usepackage[all]{xy}

% We use 2cell for 2-commutative diagrams.
\xyoption{2cell}
% Theorem environments.
%
\theoremstyle{plain}
\newtheorem{theorem}[subsection]{Theorem}
\newtheorem{proposition}[subsection]{Proposition}
\newtheorem{lemma}[subsection]{Lemma}

\theoremstyle{definition}
\newtheorem{definition}[subsection]{Definition}
\newtheorem{example}[subsection]{Example}
\newtheorem{exercise}[subsection]{Exercise}
\newtheorem{situation}[subsection]{Situation}

\theoremstyle{remark}
\newtheorem{remark}[subsection]{Remark}
\newtheorem{remarks}[subsection]{Remarks}

\numberwithin{equation}{subsection}

% Macros
%
\def\lim{\mathop{\mathrm{lim}}\nolimits}
\def\colim{\mathop{\mathrm{colim}}\nolimits}
\def\Spec{\mathop{\mathrm{Spec}}}
\def\Hom{\mathop{\mathrm{Hom}}\nolimits}
\def\Ext{\mathop{\mathrm{Ext}}\nolimits}
\def\SheafHom{\mathop{\mathcal{H}\!\mathit{om}}\nolimits}
\def\SheafExt{\mathop{\mathcal{E}\!\mathit{xt}}\nolimits}
\def\Sch{\mathit{Sch}}
\def\Mor{\mathop{\mathrm{Mor}}\nolimits}
\def\Ob{\mathop{\mathrm{Ob}}\nolimits}
\def\Sh{\mathop{\mathit{Sh}}\nolimits}
\def\NL{\mathop{N\!L}\nolimits}
\def\CH{\mathop{\mathrm{CH}}\nolimits}
\def\proetale{{pro\text{-}\acute{e}tale}}
\def\etale{{\acute{e}tale}}
\def\QCoh{\mathit{QCoh}}
\def\Ker{\mathop{\mathrm{Ker}}}
\def\Im{\mathop{\mathrm{Im}}}
\def\Coker{\mathop{\mathrm{Coker}}}
\def\Coim{\mathop{\mathrm{Coim}}}

% Boxtimes
%
\DeclareMathSymbol{\boxtimes}{\mathbin}{AMSa}{"02}

%
% Macros for moduli stacks/spaces
%
\def\QCohstack{\mathcal{QC}\!\mathit{oh}}
\def\Cohstack{\mathcal{C}\!\mathit{oh}}
\def\Spacesstack{\mathcal{S}\!\mathit{paces}}
\def\Quotfunctor{\mathrm{Quot}}
\def\Hilbfunctor{\mathrm{Hilb}}
\def\Curvesstack{\mathcal{C}\!\mathit{urves}}
\def\Polarizedstack{\mathcal{P}\!\mathit{olarized}}
\def\Complexesstack{\mathcal{C}\!\mathit{omplexes}}
% \Pic is the operator that assigns to X its picard group, usage \Pic(X)
% \Picardstack_{X/B} denotes the Picard stack of X over B
% \Picardfunctor_{X/B} denotes the Picard functor of X over B
\def\Pic{\mathop{\mathrm{Pic}}\nolimits}
\def\Picardstack{\mathcal{P}\!\mathit{ic}}
\def\Picardfunctor{\mathrm{Pic}}
\def\Deformationcategory{\mathcal{D}\!\mathit{ef}}
\usepackage{cleveref}
\crefname{originalreference}{original reference}{original references}
\Crefname{originalreference}{Original reference}{Original references}
\makeatletter
\newlabel{remark-completion-bivariant}{{0FE0}{source}{}{sourcepointer-1}{}}
\newlabel{remark-completion-bivariant@cref}{{[originalreference][0][]0FE0}{[1][0][]source}{}{}{}}
\makeatother
\crefname{originallemma}{original lemma}{original lemmas}
\Crefname{originallemma}{Original lemma}{Original lemmas}
\makeatletter
\newlabel{lemma-pullback-virtual-normal-sheaf}{{0FBJ}{source}{}{sourcepointer-2}{}}
\newlabel{lemma-pullback-virtual-normal-sheaf@cref}{{[originallemma][0][]0FBJ}{[1][0][]source}{}{}{}}
\makeatother
\crefname{originallemma}{original lemma}{original lemmas}
\Crefname{originallemma}{Original lemma}{Original lemmas}
\makeatletter
\newlabel{lemma-flat-pullback-bivariant}{{0B78}{source}{}{sourcepointer-3}{}}
\newlabel{lemma-flat-pullback-bivariant@cref}{{[originallemma][0][]0B78}{[1][0][]source}{}{}{}}
\makeatother
\crefname{originallemma}{original lemma}{original lemmas}
\Crefname{originallemma}{Original lemma}{Original lemmas}
\makeatletter
\newlabel{lemma-push-proper-bivariant}{{0EPK}{source}{}{sourcepointer-4}{}}
\newlabel{lemma-push-proper-bivariant@cref}{{[originallemma][0][]0EPK}{[1][0][]source}{}{}{}}
\makeatother
\crefname{originallemma}{original lemma}{original lemmas}
\Crefname{originallemma}{Original lemma}{Original lemmas}
\makeatletter
\newlabel{lemma-vectorbundle}{{02TY}{source}{}{sourcepointer-5}{}}
\newlabel{lemma-vectorbundle@cref}{{[originallemma][0][]02TY}{[1][0][]source}{}{}{}}
\makeatother
\crefname{originalreference}{original reference}{original references}
\Crefname{originalreference}{Original reference}{Original references}
\makeatletter
\newlabel{remark-restriction-bivariant}{{0F9Z}{source}{}{sourcepointer-6}{}}
\newlabel{remark-restriction-bivariant@cref}{{[originalreference][0][]0F9Z}{[1][0][]source}{}{}{}}
\makeatother
\crefname{originallemma}{original lemma}{original lemmas}
\Crefname{originallemma}{Original lemma}{Original lemmas}
\makeatletter
\newlabel{lemma-gysin-at-infty}{{0F9J}{source}{}{sourcepointer-7}{}}
\newlabel{lemma-gysin-at-infty@cref}{{[originallemma][0][]0F9J}{[1][0][]source}{}{}{}}
\makeatother
\crefname{originallemma}{original lemma}{original lemmas}
\Crefname{originallemma}{Original lemma}{Original lemmas}
\makeatletter
\newlabel{lemma-gysin-at-infty-independent}{{0FAQ}{source}{}{sourcepointer-8}{}}
\newlabel{lemma-gysin-at-infty-independent@cref}{{[originallemma][0][]0FAQ}{[1][0][]source}{}{}{}}
\makeatother
\crefname{originallemma}{original lemma}{original lemmas}
\Crefname{originallemma}{Original lemma}{Original lemmas}
\makeatletter
\newlabel{lemma-flat-pullback-proper-pushforward}{{02RG}{source}{}{sourcepointer-9}{}}
\newlabel{lemma-flat-pullback-proper-pushforward@cref}{{[originallemma][0][]02RG}{[1][0][]source}{}{}{}}
\makeatother
\title{The virtual normal sheaf construction is well defined as a bivariant class and is compatible with base change}
\author{The Stacks Project Authors}
\date{stacks-project:chow.tex:local-snapshot-20261012; GNU FDL 1.2 or later}
\begin{document}
\maketitle
\noindent\textit{Editorial note.} All delivered mathematical text is the author's own, copied from the pinned Stacks Project chapter file chow.tex; no statement or proof step was rewritten, reordered or removed.\par\medskip
\noindent\textit{Editorial note.} Preamble provenance: the pinned chapter file opens with an input of the shared preamble and carries no preamble text of its own, so the project's single Stacks-wide preamble staged at sources/stacks-chow/preamble.tex is excerpted verbatim as two spans, lines 14-18 (Xy-pic setup) and lines 154-217 (the native theorem declarations and mathematical macros). Nothing was added, removed or redefined; the wrapper class and the amsmath/amssymb/amsthm/graphicx base packages are supplied by the assembler.\par\medskip
\noindent\textit{Editorial note.} The counted statement and proof contain no cross reference, no citation and no included file, so the delivered body is self contained.\par\medskip
\noindent\textit{Editorial note.} Canonical location of the counted proposition: Stacks tag 0FBK (sources/stacks-chow/pinned-original-source.tex line 11912).\par\medskip
\noindent\textit{Allowed standard prerequisites (editorial).}\begin{itemize}
\item Flat pullback, proper pushforward and the class C of the Gysin map at infinity as bivariant classes. Same chapter, Lemmas labelled lemma-flat-pullback-bivariant, lemma-push-proper-bivariant and lemma-gysin-at-infty, cited at proof lines 11926-11931 and 11960-11962. Makes both sides of the defining equation bivariant classes and supplies the class C used to define c.
\item Bijectivity of the pullback along a vector bundle. Same chapter, Lemma labelled lemma-vectorbundle, cited at proof lines 11932-11934. Allows c to be defined from the equation by dividing by p\textasciicircum{}*.
\item Base change of the virtual normal sheaf and the independence of C under pushforward. Same chapter, Lemma labelled lemma-pullback-virtual-normal-sheaf, cited at proof lines 11911-11915 and 11938-11940, and Lemma labelled lemma-gysin-at-infty-independent, cited at proof lines 11972-11976. Fixes the class N' of the base changed situation and compares the two classes C cap alpha' and C' cap alpha'.
\item Restriction of bivariant classes and flat pullback of proper pushforwards. Same chapter, Remark labelled remark-restriction-bivariant, invoked in the statement, and Lemma labelled lemma-flat-pullback-proper-pushforward, cited at proof lines 11978-11982. Gives the meaning of the compatibility statement and transports the identity to the opens.
\end{itemize}\medskip
\ifdefined\hypertarget\hypertarget{sourcepointer-1}{}\fi
\noindent\textit{Source reference.} Original reference 0FE0: sources/stacks-chow/pinned-original-source.tex, source lines 6675--6675. Same-chapter pointer into the author environment remark at source line 6675; invoked at proof line None of the counted proof. It is retained only as a navigation pointer; every load-bearing step is inside the delivered span.\par\medskip
\ifdefined\hypertarget\hypertarget{sourcepointer-2}{}\fi
\noindent\textit{Source reference.} Original lemma 0FBJ: sources/stacks-chow/pinned-original-source.tex, source lines 11851--11851. Same-chapter pointer into the author environment lemma at source line 11851; invoked at proof line 11939 of the counted proof. It is retained only as a navigation pointer; every load-bearing step is inside the delivered span.\par\medskip
\ifdefined\hypertarget\hypertarget{sourcepointer-3}{}\fi
\noindent\textit{Source reference.} Original lemma 0B78: sources/stacks-chow/pinned-original-source.tex, source lines 6131--6131. Same-chapter pointer into the author environment lemma at source line 6131; invoked at proof line 11928 of the counted proof. It is retained only as a navigation pointer; every load-bearing step is inside the delivered span.\par\medskip
\ifdefined\hypertarget\hypertarget{sourcepointer-4}{}\fi
\noindent\textit{Source reference.} Original lemma 0EPK: sources/stacks-chow/pinned-original-source.tex, source lines 6167--6167. Same-chapter pointer into the author environment lemma at source line 6167; invoked at proof line 11929 of the counted proof. It is retained only as a navigation pointer; every load-bearing step is inside the delivered span.\par\medskip
\ifdefined\hypertarget\hypertarget{sourcepointer-5}{}\fi
\noindent\textit{Source reference.} Original lemma 02TY: sources/stacks-chow/pinned-original-source.tex, source lines 7062--7062. Same-chapter pointer into the author environment lemma at source line 7062; invoked at proof line 11935 of the counted proof. It is retained only as a navigation pointer; every load-bearing step is inside the delivered span.\par\medskip
\ifdefined\hypertarget\hypertarget{sourcepointer-6}{}\fi
\noindent\textit{Source reference.} Original reference 0F9Z: sources/stacks-chow/pinned-original-source.tex, source lines 6184--6184. Same-chapter pointer into the author environment remark at source line 6184; invoked at proof line 11942 of the counted proof. It is retained only as a navigation pointer; every load-bearing step is inside the delivered span.\par\medskip
\ifdefined\hypertarget\hypertarget{sourcepointer-7}{}\fi
\noindent\textit{Source reference.} Original lemma 0F9J: sources/stacks-chow/pinned-original-source.tex, source lines 10026--10026. Same-chapter pointer into the author environment lemma at source line 10026; invoked at proof line 11968 of the counted proof. It is retained only as a navigation pointer; every load-bearing step is inside the delivered span.\par\medskip
\ifdefined\hypertarget\hypertarget{sourcepointer-8}{}\fi
\noindent\textit{Source reference.} Original lemma 0FAQ: sources/stacks-chow/pinned-original-source.tex, source lines 10108--10108. Same-chapter pointer into the author environment lemma at source line 10108; invoked at proof line 11971 of the counted proof. It is retained only as a navigation pointer; every load-bearing step is inside the delivered span.\par\medskip
\ifdefined\hypertarget\hypertarget{sourcepointer-9}{}\fi
\noindent\textit{Source reference.} Original lemma 02RG: sources/stacks-chow/pinned-original-source.tex, source lines 2401--2401. Same-chapter pointer into the author environment lemma at source line 2401; invoked at proof line 11980 of the counted proof. It is retained only as a navigation pointer; every load-bearing step is inside the delivered span.\par\medskip
\begin{lemma}
\label{lemma-construction-gysin}
The construction above defines a bivariant class\footnote{The
notation $A^*(Z \to X)^\wedge$ is discussed in
Remark \ref{remark-completion-bivariant}.
If $X$ is quasi-compact, then $A^*(Z \to X)^\wedge = A^*(Z \to X)$.}
$$
c(Z \to X, \mathcal{N}) \in A^*(Z \to X)^\wedge
$$
and moreover the construction is compatible with base change
as in Lemma \ref{lemma-pullback-virtual-normal-sheaf}.
If $\mathcal{N}$ has constant rank $r$, then
$c(Z \to X, \mathcal{N}) \in A^r(Z \to X)$.
\end{lemma}

\begin{proof}
Since both $i_* \circ j^* \circ C$ and $p^*$ are bivariant classes
(see Lemmas \ref{lemma-flat-pullback-bivariant} and
\ref{lemma-push-proper-bivariant}) we can use the equation
$$
i_* \circ j^* \circ C  = p^* \circ c(Z \to X, \mathcal{N})
$$
(suitably interpreted) to define $c(Z \to X, \mathcal{N})$
as a bivariant class. This works because $p^*$ is always
bijective on chow groups by Lemma \ref{lemma-vectorbundle}.

\medskip\noindent
Let $X' \to X$, $Z' \to X'$, and $\mathcal{N}'$ be as in
Lemma \ref{lemma-pullback-virtual-normal-sheaf}. Write
$c = c(Z \to X, \mathcal{N})$ and $c' = c(Z' \to X', \mathcal{N}')$.
The second statement of the lemma means that $c'$ is the restriction of $c$
as in Remark \ref{remark-restriction-bivariant}. Since we claim this
is true for all $X'/X$ locally of finite type, a formal argument
shows that it suffices to check that $c' \cap \alpha' = c \cap \alpha'$
for $\alpha' \in \CH_k(X')$.
To see this, note that we have a commutative diagram
$$
\xymatrix{
C_{Z'}X' \ar[d] \ar[r] &
W'_\infty \ar[d] \ar[r] &
W' \ar[d] \ar[r] &
\mathbf{P}^1_{X'} \ar[d] \\
C_ZX \ar[r] &
W_\infty \ar[r] &
W \ar[r] &
\mathbf{P}^1_X
}
$$
which induces closed immersions:
$$
W' \to W \times_{\mathbf{P}^1_X} \mathbf{P}^1_{X'},\quad
W'_\infty \to W_\infty \times_X X',\quad
C_{Z'}X' \to C_ZX \times_Z Z'
$$
To get $c \cap \alpha'$ we use the class $C \cap \alpha'$
defined using the morphism
$W \times_{\mathbf{P}^1_X} \mathbf{P}^1_{X'} \to \mathbf{P}^1_{X'}$
in Lemma \ref{lemma-gysin-at-infty}.
To get $c' \cap \alpha'$ on the other hand, we use the class
$C' \cap \alpha'$ defined using the morphism $W' \to \mathbf{P}^1_{X'}$.
By Lemma \ref{lemma-gysin-at-infty-independent} the pushforward of
$C' \cap \alpha'$ by the closed immersion
$W'_\infty \to (W \times_{\mathbf{P}^1_X} \mathbf{P}^1_{X'})_\infty$,
is equal to $C \cap \alpha'$. Hence the same is true for the pullbacks
to the opens
$$
C_{Z'}X' \subset W'_\infty,\quad
C_ZX \times_Z Z' \subset (W \times_{\mathbf{P}^1_X} \mathbf{P}^1_{X'})_\infty
$$
by Lemma \ref{lemma-flat-pullback-proper-pushforward}.
Since we have a commutative diagram
$$
\xymatrix{
C_{Z'} X' \ar[d] \ar[r] & N' \ar@{=}[d] \\
C_ZX \times_Z Z' \ar[r] & N \times_Z Z'
}
$$
these classes pushforward to the same class on $N'$ which
proves that we obtain the same element $c \cap \alpha' = c' \cap \alpha'$
in $\CH_*(Z')$.
\end{proof}

\end{document}
